\chapter{Conclusion and Discussion}
\label{chap:conclusion-discussion}

\section{Challenges and Solutions}
\label{section:challenges}

\subsection{Finding an Effective Low-Noise Poisoning Method}

The most significant challenge in this project was
identifying an adversarial perturbation method that
could cause a video action recognition model to
misclassify while keeping the noise imperceptibly small.
Most published research in this area focuses on
high-visibility perturbations or physical-world attacks,
such as adversarial patches or printed eyeglass frames,
which are not suitable for the use case of protecting
user-uploaded video content.

The team surveyed multiple attack methods from the
literature, including universal adversarial perturbations,
black-box transfer attacks, and feature-space poisoning
approaches, before settling on the sparse adversarial
spatial-transform approach from Mu et al. (BMVC 2021).
Even after selecting this method, adapting it to work
reliably on SlowFast R101 required extensive
experimentation with the noise budget, SSIM constraint,
Bayesian Optimisation parameters, and the codec boost
factor to compensate for compression attenuation.

\subsection{GPU Compatibility with Bayesian Frame Selection}

The Bayesian frame selection component runs a short
inner optimisation loop for each candidate mask
evaluation. On CPU this is manageable for a small
number of candidates, but the full BO procedure with
50 iterations is computationally intensive. Testing
on Apple Silicon MacBooks revealed that the MPS
backend does not fully support the tensor operations
used in the spatial transformer and the GP kernel
computations, causing the pipeline to fall back to
CPU. A single BO pass on CPU without GPU acceleration
was observed to take approximately one hour per video,
which made iterative development impractical.

This was partially addressed by reducing the number of
BO iterations during development and running full
evaluations on CUDA-enabled hardware. The final
pipeline defaults to CUDA when available and falls
back to CPU otherwise, with a note to users that
processing time on CPU may be significantly longer.

\subsection{Dependency on a Large Pretrained Model}

SlowFast R101 is a large model that must be downloaded
from the PyTorchVideo hub on first use. This creates
a dependency on an internet connection during the
initial setup and adds approximately 200 MB to the
application footprint. The model is cached after the
first download, but users on slow connections or
restricted networks may encounter difficulties. This
was mitigated by adding a loading progress indicator
in the application and providing clear error messages
when the model cannot be fetched.

\subsection{Background in AI and Data Poisoning}

At the start of the project, both team members had
limited prior experience with adversarial machine
learning and video model architectures. To build the
necessary background, the team participated in a
series of AI workshop labs advised by the project
supervisor, beginning in January 2025. Topics covered
included binary image classification, the stAdv
spatial-transform attack, text-to-video model training,
and quantifiable metrics for evaluating model
degradation under adversarial inputs. This foundational
work was essential for understanding the limitations of
existing approaches and making informed design decisions
throughout the project.

\section{Limitations}
\label{section:limitations}

\subsection{GPU Requirement for Practical Use}

The full VidScratch pipeline with Bayesian Optimisation
is not practical to run on CPU for most users. On a
modern laptop without a CUDA-capable GPU, processing
a single short video can take over an hour. This
limits the accessibility of the application for
users who do not have access to a discrete GPU,
which contradicts the project's goal of providing an
accessible tool for everyday users.

\subsection{Limited Impact on Training-Time Attacks}

The training-time poisoning experiment demonstrated
that the perturbations produced by VidScratch have
only a modest effect on a model trained on poisoned
data. The adversarial noise is optimised to fool a
400-class SlowFast R101 model, but when the downstream
model is trained on a small two-class dataset the
gradient signal from the perturbation is not well
aligned with the training objective. The result is
a subtle increase in test loss rather than a clear
degradation in model performance. For the training-time
attack to be more effective, the perturbation would
need to be generated with knowledge of the downstream
training setup, which is not available in a practical
user-facing tool.

\subsection{Fixed Target Model}

The current implementation targets SlowFast R101
exclusively. Users who want to test robustness
against a different action recognition model, or
who encounter a platform that uses a different
architecture, cannot use VidScratch directly without
modifying the codebase. The pipeline is designed to
be extensible, but swapping the target model requires
technical knowledge of the Python backend.

\subsection{Short Video Constraint}

The pipeline is designed for videos of up to one
minute in length, which covers most short-form
content but excludes longer recordings. Extending
support to longer videos would require changes to
the frame sampling and memory management logic,
as processing all frames of a long video at full
resolution would exceed the memory available on
most consumer hardware.

\section{Future Work}
\label{section:future-work}

\subsection{Support for Transformer-Based Video Models}

The most natural extension of this project is to
target more recent video understanding architectures,
particularly transformer-based models such as
VideoMAE, TimeSformer, and Video Swin Transformer.
These models have largely replaced CNN-based
architectures such as SlowFast on standard benchmarks.
Generating adversarial perturbations that transfer
across model families or that are effective against
attention-based temporal reasoning specifically
remains an open and technically interesting problem.

\subsection{Improved Training-Time Effectiveness}

Future work could explore poisoning strategies that
are more directly aligned with the downstream training
objective. One direction is to generate perturbations
that target the feature representations learned by the
model rather than its final classification output.
Targeting intermediate feature layers tends to
produce more transferable and training-disruptive
perturbations, and may remain effective even when
the downstream model differs from the attack model.

\subsection{GPU-Efficient Frame Selection}

Replacing the current Gaussian-Process-based Bayesian
Optimisation with a lighter alternative, such as a
gradient-based saliency score or a simple frame-level
loss ranking, could substantially reduce processing
time without significantly reducing attack effectiveness.
This would make the pipeline usable on CPU in a
practical amount of time and improve accessibility
for users without dedicated GPU hardware.

\subsection{Broader Video Format Support}

The current version supports MP4 input and output only.
Future versions could extend support to additional
formats such as AVI, MKV, and MOV, and could also
support frame-rate-aware perturbation to handle
high-frame-rate content more consistently.

\section{Privacy and Data Usage}
\label{section:privacy}

Vid-Scratch processes all video data locally on the
user's machine. No video content, frame data, or
personal information is transmitted to any external
server at any point during the poisoning pipeline.
The application does not collect usage data or
telemetry. Adversarial outputs produced by the
pipeline are saved to a user-specified local folder
and are not shared or made accessible to any third
party.

The project did not use any copyrighted material
without permission during development. Videos used
for evaluation were drawn from the publicly available
Kinetics-400 and UCF-101 datasets under their
respective research-use licenses. Any user-supplied
video content is used solely for the purpose of
generating the adversarial output requested by the
user and is not retained by the application.

\section{Conclusion}
\label{section:conclusion}

This project presented Vid-Scratch, a desktop
application that generates adversarially perturbed
video clips designed to cause a pre-trained action
recognition model to misclassify the action shown
in the input video, while preserving high visual
similarity to the original content. The system
targets SlowFast R101 pretrained on Kinetics-400
and implements a two-stage pipeline consisting of
Bayesian Optimisation for sparse frame selection
followed by joint optimisation of a spatial
transformation field and additive noise using Adam.

The inference-time attack achieved a verified
success rate of 100\% across 10 Kinetics-400 test
videos with an average SSIM of 0.9501, demonstrating
that the approach can reliably cause misclassification
while keeping the perturbation imperceptible to
casual inspection. The training-time experiment
showed a modest but observable effect on clean
test loss when poisoned samples were mixed into
the training data, though the limited dataset
size and the mismatch between the 400-class attack
model and the 2-class training setup reduce the
strength of this result.

Beyond the technical results, the project also
produced an accessible desktop application that
allows non-expert users to apply adversarial
video poisoning without requiring knowledge of
the underlying machine learning methods. In this
respect, Vid-Scratch bridges adversarial machine
learning research and a more practical application
setting, and serves as a demonstration of how
strong benchmark performance in action recognition
does not necessarily imply robustness under
carefully constructed adversarial inputs.

Several limitations remain, most notably the
requirement for a CUDA-capable GPU for practical
use and the limited effectiveness of the
training-time attack under the current noise
budget. These point to clear directions for
future work, including support for
transformer-based video models, GPU-efficient
frame selection, and more targeted training-time
poisoning strategies. With these improvements,
a future version of the system could provide a
more complete and practically useful tool for
evaluating and demonstrating the adversarial
vulnerability of video understanding models.
